考虑这句话：``班上所有戴眼镜的同学数学都及格了。''要反驳它，你必须找到一个戴眼镜但数学不及格的同学；要证明它，你得检查每一个戴眼镜的同学。反之，``班上有一个同学数学满分''——证明它只需要举出一个人，反驳它却要检查所有人。一个量词的差异，彻底改变了证明和反例的形态。

一句话即使语法通顺，也可能没有确定真假，或者把条件方向、讨论范围和否定对象说得含混。逻辑的作用不是把日常语言换成符号，而是把一句话到底断定了什么变成可以逐项检查的结构。

\subsection{怎么读这一章}

三篇构成一条核心主线，建议第一次按顺序阅读：

\begin{enumerate}
\tightlist
\item \hyperref[sec:01-Mathematical-Language/01-Logic/01]{``命题与逻辑连接词''}：先判断一句话能否谈真假，再看复合命题怎样由局部真假组成；
\item \hyperref[sec:01-Mathematical-Language/01-Logic/02]{``蕴含与充要条件''}：把``若''、``仅当''和``当且仅当''转成有方向的逻辑关系，并学会构造精准反例；
\item \hyperref[sec:01-Mathematical-Language/01-Logic/03]{``量词与否定''}：理解``所有''与``存在''的证据要求，以及量词顺序为什么会改变命题。
\end{enumerate}

若眼前问题是读不懂定理条件，可以直接从第二篇进入；若困难是不会否定一句复杂陈述，可以直接读第三篇，再回补前两篇。

\subsection{读完后应发生什么变化}

目标不是背下真值表，而是形成一套可迁移的阅读动作：看到任意数学陈述，能够指出讨论对象的范围、区分假设与结论、判断怎样的证据足以证明它，并提前说出反例必须具有的形状。后面的证明、极限、概率和统计推断都会反复调用这些动作。
